\documentclass[11pt,a4paper]{article}
\usepackage[margin=1in]{geometry}
\usepackage{kotex}
\usepackage{amsmath,amssymb,bm}
\usepackage{booktabs}
\usepackage{hyperref}
\usepackage{enumitem}
\usepackage{array}
\usepackage{xcolor}
\usepackage{listings}
\usepackage{tikz}
\usetikzlibrary{arrows.meta,positioning,calc,fit,shapes.geometric,matrix,backgrounds}

\lstset{
  language=Python,
  basicstyle=\ttfamily\footnotesize,
  commentstyle=\color{green!40!black},
  keywordstyle=\color{blue!70!black},
  stringstyle=\color{orange!60!black},
  showstringspaces=false,
  columns=fullflexible,
  keepspaces=true,
  frame=single,
  breaklines=true
}

\title{Hands-On Graph Neural Networks Using Python\\Chapter 10 상세 정리}
\author{A7610 그래프신경망특론 / 조남언}
\date{\today}

\begin{document}
\maketitle
\tableofcontents
\newpage

\section{개요: 링크 예측이라는 과제}
링크 예측(link prediction)은 그래프 $G=(V,E)$에서 \textbf{관측되지 않은} 노드 쌍 $(u,v)\notin E$가
실제로 연결되어 있을(또는 미래에 연결될) 확률을 추정하는 문제다. 추천 시스템의 사용자-아이템 연결,
지식 그래프의 누락 사실 보완, 단백질-단백질 상호작용 예측 등 응용 범위가 매우 넓다.

\subsection{전통적 휴리스틱}
도메인 지식에 기반한 점수 함수들이 오랫동안 활용되었다.
\begin{itemize}[leftmargin=*]
  \item Common Neighbors: $s(u,v) = |\Gamma(u)\cap\Gamma(v)|$
  \item Jaccard: $|\Gamma(u)\cap\Gamma(v)|/|\Gamma(u)\cup\Gamma(v)|$
  \item Adamic--Adar: $\sum_{w\in\Gamma(u)\cap\Gamma(v)} 1/\log|\Gamma(w)|$
  \item Katz: 모든 경로 길이를 감쇠 계수로 합산
\end{itemize}
이들은 \emph{enclosing subgraph}(두 노드를 둘러싼 작은 구조)의 함수라는 공통점이 있다.
SEAL은 바로 이 점에서 출발해 GNN으로 \emph{그 함수 자체}를 학습한다.

\subsection{데이터 분할의 특수성}
링크 예측은 노드 분류와 달리 \textbf{엣지} 단위로 train/val/test를 나눈다.
PyG의 \texttt{RandomLinkSplit}은 그래프를 잘라낸 후 양성/음성 엣지 인덱스를 함께 제공한다.
평가 지표로는 AUC와 AP(평균 정밀도)가 표준이다.

\section{Graph Autoencoder (GAE)}
\subsection{모델 구조}
GAE는 단순한 인코더-디코더 구조다.
\begin{align*}
  \text{인코더}: & \quad \mathbf{Z} = \mathrm{GCN}(X,A) \in \mathbb{R}^{N\times d} \\
  \text{디코더}: & \quad \hat{A} = \sigma(\mathbf{Z}\mathbf{Z}^\top)
\end{align*}
디코더가 \emph{내적}이라는 점이 핵심이다. 이는 ``비슷한 노드끼리 임베딩 공간에서 가깝다''라는 단순한
유사도 가설에 기반하며, 모든 노드 쌍에 대한 점수를 행렬곱 한 번으로 얻을 수 있어 추론이 빠르다.

\subsection{학습 손실}
양성 엣지 $E^+$와 무작위 음성 엣지 $E^-$에 대한 이진 교차엔트로피:
\begin{equation*}
  \mathcal{L}_{\text{recon}} =
   -\!\!\sum_{(u,v)\in E^{+}} \log \hat{A}_{uv}
   -\!\!\sum_{(u,v)\in E^{-}} \log(1-\hat{A}_{uv}).
\end{equation*}
클래스 불균형이 심하므로 양/음 같은 수로 샘플링한다. 한계는 명백하다: 임베딩이 결정적이라
모델의 불확실성을 표현할 수 없다.

\section{Variational Graph Autoencoder (VGAE)}
\subsection{확률적 인코더}
VGAE는 잠재변수 $\mathbf{z}_i$를 가우시안 분포로 본다.
\begin{equation*}
  q(\mathbf{z}_i\mid X,A) = \mathcal{N}(\bm{\mu}_i, \mathrm{diag}(\bm{\sigma}_i^2)),
  \quad
  \bm{\mu} = \mathrm{GCN}_{\mu}(X,A),\ \log\bm{\sigma} = \mathrm{GCN}_{\sigma}(X,A).
\end{equation*}
재매개화 트릭으로 학습 가능한 샘플을 뽑는다: $\mathbf{z} = \bm{\mu} + \bm{\sigma}\odot\bm{\epsilon}$, $\bm{\epsilon}\sim\mathcal{N}(0,I)$.
디코더는 GAE와 동일하다.

\subsection{ELBO와 KL 정규화}
변분 추론의 ELBO를 최소화 형태로 쓰면:
\begin{equation*}
  \mathcal{L}_{\text{VGAE}} = \mathcal{L}_{\text{recon}}
    + \tfrac{1}{N}\,\mathrm{KL}\bigl[q(Z\mid X,A)\,\|\,p(Z)\bigr],
\end{equation*}
사전분포 $p(Z)=\mathcal{N}(0,I)$이며 KL은 정규화 효과를 준다. $1/N$은 노드 수에 비례한 KL 누적을 노드당 평균으로
환원하기 위한 가중치이며, 책의 구현에서 사용된 관용적 형태다.

\subsection{PyG 구현 핵심}
\begin{lstlisting}
from torch_geometric.nn import GCNConv, VGAE

class Encoder(torch.nn.Module):
    def __init__(self, dim_in, dim_out):
        super().__init__()
        self.conv1       = GCNConv(dim_in, 2 * dim_out)
        self.conv_mu     = GCNConv(2 * dim_out, dim_out)
        self.conv_logstd = GCNConv(2 * dim_out, dim_out)
    def forward(self, x, edge_index):
        x = self.conv1(x, edge_index).relu()
        return (self.conv_mu(x, edge_index),
                self.conv_logstd(x, edge_index))

model = VGAE(Encoder(dataset.num_features, 16)).to(device)
\end{lstlisting}
\texttt{VGAE} 래퍼가 \texttt{recon\_loss}, \texttt{kl\_loss}, \texttt{test}를 모두 제공하므로
학습 코드는 매우 짧다.

\subsection{Cora 실험 결과}
- Encoder: GCN(1433→32) → GCN(32→16)$\times$2.\\
- 300 epoch, Adam (lr=0.01).\\
- 결과: 테스트 AUC $\approx$ 0.88, AP $\approx$ 0.90.\\
학습이 끝나면 임의의 노드 쌍에 대한 점수를 $\hat{A}=\sigma(ZZ^\top)$로 한꺼번에 얻을 수 있다.

\section{SEAL: 서브그래프 분류로서의 링크 예측}
\subsection{문제 재정의}
SEAL은 링크 예측을 \textbf{이진 그래프 분류}로 환원한다. 후보 링크 $(u,v)$의
$k$-hop \emph{enclosing subgraph}를 추출해, ``이 서브그래프는 두 노드가 연결된 상태인가?''
를 GNN으로 판단한다.

\subsection{Enclosing Subgraph 추출과 타깃 누설 방지}
\texttt{k\_hop\_subgraph}로 후보 두 노드의 2-hop 이웃 합집합을 얻은 뒤, 학습 시에 모델이
정답을 보지 않도록 \textbf{타깃 엣지 자체를 서브그래프에서 제거}한다.
\begin{lstlisting}
sub_nodes, sub_edge_index, mapping, _ = k_hop_subgraph(
    [src, dst], 2, dataset.edge_index, relabel_nodes=True)
src, dst = mapping.tolist()
mask1 = (sub_edge_index[0] != src) | (sub_edge_index[1] != dst)
mask2 = (sub_edge_index[0] != dst) | (sub_edge_index[1] != src)
sub_edge_index = sub_edge_index[:, mask1 & mask2]
\end{lstlisting}
이 처리는 양성 샘플뿐 아니라 음성 샘플에도 동일하게 적용해야 분포가 일치한다.

\subsection{DRNL (Double-Radius Node Labeling)}
GNN은 서브그래프 안의 어떤 노드가 ``타깃''인지 모른다. 따라서 거리 정보를 \emph{명시적으로 노드 라벨로 부여}한다.
타깃이 아닌 노드 $x$에 대해 $d_x=d(x,u)$, $d_y=d(x,v)$라 하자.
$d=(d_x+d_y)/2$, $r=(d_x+d_y) \bmod 2$일 때 DRNL 라벨은
\begin{equation*}
  z(x) = 1 + \min(d_x, d_y) + d\cdot(d+r-1),
\end{equation*}
이며 타깃 두 노드는 $z=1$로 고정, 도달 불가 노드는 $z=0$이다. 두 거리를 단일 정수로 인코딩하여
원-핫 표현이 깔끔하며, 거리 기반 휴리스틱들의 표현력을 포괄한다.

\textbf{거리 계산의 한 가지 트릭}: $d_x$를 계산할 때는 그래프에서 $v$를 잠시 제외하고,
$d_y$를 계산할 때는 $u$를 제외한다. 그렇게 해야 두 거리가 ``상대방을 거쳐가는 경로''로 줄어드는 누설을 막을 수 있다.

\subsection{DGCNN}
SEAL의 분류기는 DGCNN이다.
\begin{enumerate}[leftmargin=*]
  \item 4개의 GCN 레이어(32→32→32→1, \texttt{tanh} 활성)의 출력을 \textbf{레이어별 concat}.
  \item \textbf{SortPooling}: 마지막 채널 값으로 노드를 정렬해 상위 $k=30$개만 남김.
        그래프 크기가 달라도 \emph{고정 길이} 표현이 된다.
  \item Conv1D + MaxPool로 정렬된 시퀀스를 처리(텍스트 분류와 유사).
  \item Linear + Sigmoid로 링크 존재 확률 출력.
\end{enumerate}

\subsection{Cora 실험 결과}
- 30 epoch, Adam (lr=$10^{-4}$), BCEWithLogitsLoss.\\
- 결과: 테스트 AUC $\approx$ 0.89, AP $\approx$ 0.90. VGAE와 동급이거나 약간 우위.\\
표현력은 좋지만 \textbf{쌍마다 서브그래프 추출/전처리}가 필요해 비용이 크다.

\section{비교와 시사점}
\begin{center}
\begin{tabular}{p{0.20\linewidth} p{0.35\linewidth} p{0.35\linewidth}}
\toprule
 & \textbf{VGAE} & \textbf{SEAL} \\
\midrule
관점     & 그래프 전체 임베딩 & 후보 링크별 서브그래프 분류 \\
입력 단위 & 그래프 1개 & 후보 링크 1개 \\
디코더   & 내적 $\sigma(z_u^\top z_v)$ & DGCNN \\
노드 라벨 & 원본 feature & 원본 feature + DRNL one-hot \\
학습 비용 & 작음 & 큼 (전처리 비싼) \\
추론 비용 & 매우 작음 ($O(d)$/쌍) & 큼 (쌍마다 서브그래프) \\
표현력   & 대칭/저차원 한계 & 휴리스틱 일반화, 풍부 \\
\bottomrule
\end{tabular}
\end{center}

\paragraph{언제 무엇을 쓰나}
\begin{itemize}[leftmargin=*]
  \item 노드 수가 크고 후보 쌍이 많다면 VGAE처럼 \emph{한 번 학습-즉시 점수} 모델이 유리.
  \item 적은 양이라도 정밀한 후보 평가가 필요하면 SEAL이 우수.
  \item 두 접근은 결합 가능 — VGAE로 후보 풀을 거른 뒤 SEAL로 재정렬하는 식의 캐스케이드도 자연스럽다.
\end{itemize}

\paragraph{이번 장의 의의}
GAE/VGAE는 \emph{표현 학습}의 관점, SEAL은 \emph{문제 환원}의 관점이다.
같은 과제를 어떤 단위로 보느냐가 모델 구조와 비용 특성을 전혀 다르게 만든다는 점이 핵심.

\end{document}
